% ============================================================
%  problem_formulation.tex  —  §3
%  Updated: 29/Sep/2026
% ============================================================

\section{Problem formulation}
\label{sec:problem}

\subsection{Candidate sites and the separation constraint}

Let $S=\{1,\dots,N\}$ index candidate sites, each described by $m=11$
geospatial criteria.
A \emph{portfolio} is a subset $t\subseteq S$ with $|t|=\kk$.
Portfolios must satisfy a hard pairwise separation constraint
\begin{equation}
  \label{eq:rsep}
  d(u,v) \;\ge\; d_{\min}=\dmin\;\text{km}
  \qquad \forall\, u,v\in t,\; u\neq v,
\end{equation}
where $d(\cdot,\cdot)$ is the great-circle distance.
Constraint~\eqref{eq:rsep} is the $r$-separation condition of
\citet{Erkut1996}, whose network antecedent is \citet{MoonChaudhry1984}.
We emphasise that dispersion enters here as a \emph{constraint}, not an
objective: the $p$-dispersion literature \citep{Kuby1987,Marti2022}
maximises the minimum pairwise distance, whereas~\eqref{eq:rsep} merely
renders non-compliant portfolios infeasible.
The feasible set is
\begin{equation}
  \label{eq:T}
  \Tset = \bigl\{\,t\subseteq S \;:\; |t|=\kk,\ t \text{ satisfies }
    \eqref{eq:rsep}\,\bigr\},
  \qquad |\Tset| = \Ntrios.
\end{equation}

\subsection{Three partial value functions}

The $m$ criteria are partitioned into three groups reflecting the
decision dimensions of landfill siting: siting and construction cost
($G_1$), environmental impact ($G_2$), and social safety ($G_3$)
(Table~\ref{tab:criteria}).
Let $c_j(i)\in[0,1]$ be the partial value of site~$i$ on criterion~$j$,
oriented so that larger is better, and let $w_j$ be the criterion weight.
The portfolio objectives are additive within group and across sites:
\begin{align}
  F_1(t) &= \sum_{i\in t}\;\sum_{j\in G_1} w_j\,c_j(i),
    && \text{(cost, minimise)} \label{eq:F1} \\
  F_2(t) &= \sum_{i\in t}\;\sum_{j\in G_2} w_j\,c_j(i),
    && \text{(impact, minimise)} \label{eq:F2} \\
  F_3(t) &= \sum_{i\in t}\;\sum_{j\in G_3} w_j\,c_j(i),
    && \text{(safety, maximise)} \label{eq:F3}
\end{align}
and we write $\mathbf{F}(t;w)=(F_1,F_2,F_3)$ to make the weight
dependence explicit, since it is central to
Section~\ref{sec:design}.

\begin{remark}[what the benchmark is, and is not]
\label{rem:valuespace}
$\mathbf{F}$ is a three-dimensional partial value space, not a vector of
physically independent objectives: each component is an additive weighted
aggregate of normalised criteria.
This distinguishes our reference set from portfolio decision analysis
approaches that compute non-dominated portfolios in a scalar value space
\citep{Liesio2007,Liesio2008,Yang2015}, and from bi-objective mixed-integer
landfill models whose objectives are physical quantities
\citep{EiseltMarianov2014}.
We keep the three group-level value functions unaggregated and ask how
much of the resulting trade-off surface is lost when a ranking method
collapses them into a single score before the portfolio is assembled.
\end{remark}

\subsection{Reference set and the optimality gap}

Portfolio $t$ \emph{dominates} $t'$ when $F_1(t)\le F_1(t')$,
$F_2(t)\le F_2(t')$, $F_3(t)\ge F_3(t')$ with at least one strict
inequality.
Because $|\Tset|$ is finite and enumerable, the non-dominated set
\begin{equation}
  \label{eq:front}
  \Fset(w) = \bigl\{\,t\in\Tset \;:\; \nexists\,t'\in\Tset
    \text{ dominating } t\,\bigr\}
\end{equation}
is obtained \emph{exactly}, with no metaheuristic approximation.
All dominance tests use a tolerance $\varepsilon=10^{-9}$.

For dominated portfolios we report a scale-free shortfall.
With $R_\ell = \max_{t\in\Tset}F_\ell - \min_{t\in\Tset}F_\ell$,
\begin{equation}
  \label{eq:gap}
  \mathrm{gap}(t) = \max\Bigl(0,\;\max_{t'\in\Tset}\;\min\Bigl\{
    \tfrac{F_1(t)-F_1(t')}{R_1},\;
    \tfrac{F_2(t)-F_2(t')}{R_2},\;
    \tfrac{F_3(t')-F_3(t)}{R_3}\Bigr\}\Bigr),
\end{equation}
so $\mathrm{gap}(t)=0$ iff $t\in\Fset$.
Normalising by the attainable range rather than the ideal--nadir distance
keeps the measure interpretable as a fraction of the achievable spread;
the ideal--nadir variant is reported in the Supplementary Material.
